\documentclass[10pt,twoside,openany]{book}
\usepackage[a5paper,inner=18mm,outer=14mm,top=15mm,bottom=17mm]{geometry}
\usepackage{fontspec}
\setmainfont{texgyrepagella}[Extension=.otf,UprightFont=*-regular,BoldFont=*-bold,ItalicFont=*-italic,BoldItalicFont=*-bolditalic,Ligatures=TeX]
\usepackage{xcolor}\usepackage{fancyhdr}\usepackage{tcolorbox}\usepackage{setspace}\usepackage{enumitem}\usepackage{booktabs}\usepackage{tabularx}
\tcbuselibrary{skins,breakable}
\definecolor{accent}{RGB}{154,52,18}\definecolor{gris}{gray}{0.45}\definecolor{creme}{RGB}{252,246,238}\definecolor{vert}{RGB}{31,102,70}\definecolor{bleu}{RGB}{30,64,120}
\setstretch{1.02}\setlength{\parindent}{0pt}\setlength{\parskip}{0.45em}\widowpenalty=10000\clubpenalty=10000\emergencystretch=2em
\setlist{nosep,leftmargin=1.4em,itemsep=0.25em}
\pagestyle{fancy}\fancyhf{}\renewcommand{\headrulewidth}{0pt}
\fancyhead[LE]{\small\itshape Tu n’es pas nul — le guide de poche}\fancyhead[RO]{\small\itshape\rightmark}\fancyfoot[C]{\small\thepage}
\fancypagestyle{plain}{\fancyhf{}\fancyfoot[C]{\small\thepage}\renewcommand{\headrulewidth}{0pt}}
\newenvironment{aretenir}[1][À retenir]{\par\begin{tcolorbox}[enhanced,colback=creme,colframe=accent,boxrule=0pt,leftrule=3pt,arc=0pt,left=8pt,right=8pt,top=5pt,bottom=5pt,breakable]{\small\bfseries\color{accent}\MakeUppercase{#1}}\par\smallskip}{\end{tcolorbox}\par}
\newenvironment{cesoir}[1][Ce soir, essaie ceci]{\par\begin{tcolorbox}[enhanced,colback=white,colframe=vert,boxrule=0.6pt,arc=2pt,left=8pt,right=8pt,top=5pt,bottom=5pt,breakable]{\small\bfseries\color{vert}\MakeUppercase{#1}}\par\smallskip\small}{\end{tcolorbox}\par}
\newenvironment{prudence}[1][Ce que dit la recherche]{\par\begin{tcolorbox}[enhanced,colback=black!4,colframe=bleu,boxrule=0pt,leftrule=3pt,arc=0pt,left=8pt,right=8pt,top=5pt,bottom=5pt,breakable]{\small\bfseries\color{bleu}\MakeUppercase{#1}}\par\smallskip\small}{\end{tcolorbox}\par}
\newcommand{\fiche}[3]{\clearpage\markright{#2}\thispagestyle{plain}%
 {\noindent\color{accent}\fontsize{40}{40}\selectfont\bfseries #1}\hspace{0.6em}{\Large\bfseries #2\par}\vspace{0.1em}{\small\color{gris}\itshape #3\par}\vspace{0.5em}}
\newcommand{\partie}[2]{\clearpage\thispagestyle{empty}\vspace*{\fill}\begin{center}{\large\color{gris}Partie #1}\\[1em]{\Huge\bfseries #2}\end{center}\vspace*{\fill}\clearpage}
\newcommand{\titre}[1]{\clearpage\markright{#1}\thispagestyle{plain}{\noindent\LARGE\bfseries #1\par}\vspace{0.8em}}
\begin{document}
\frontmatter
\thispagestyle{empty}
\begin{center}\vspace*{2.2cm}
{\Huge\bfseries Tu n’es pas nul\par}\vspace{0.4cm}
{\large\itshape Le guide de poche\par}\vspace{1.2cm}
{\large L’essentiel du livre en quatorze fiches,\\pour ceux qui n’ont pas le temps\\et pour ceux qui veulent un avant-goût.\par}\vspace{2cm}
{\Large Ludovic A.\par}\vfill{\small livret associé au livre — version de travail\par}\end{center}
\clearpage
\thispagestyle{empty}\vspace*{\fill}\begin{center}\itshape On a écrit un mot sur toi.\\Ce n’était pas ton nom.\end{center}\vspace*{\fill}
\mainmatter
\titre{Avant de commencer}
Ce livret est fait pour deux sortes de lecteurs. Ceux qui n’ont pas le temps de lire un livre entier, et qui veulent tout de suite quelque chose à faire. Ceux qui se demandent si le livre est pour eux, et qui veulent en goûter un morceau.

Le livre « Tu n’es pas nul » raconte des histoires : un professeur de Bafoussam qui calcule des rangs à la lueur d’une lampe qui s’éteint, une élève de Douala qui n’ose pas lever la main, un garçon de Ngaoundéré qui répare des ventilateurs. Ce livret garde seulement ce qui sert : \textbf{l’idée, la raison, et un geste à faire ce soir}. Les histoires, les nuances et les sources sont dans le livre.

\begin{aretenir}[L’idée du livre en une phrase]
Une note est un résultat sur une copie, un jour donné. Elle n’est pas ton nom. Derrière presque chaque « je suis nul » se cache quelque chose de précis, que l’on peut nommer, et que l’on peut travailler.
\end{aretenir}

\textbf{Comment s’en servir.}
\begin{itemize}
\item Prends un cahier de brouillon. Il te servira pour tous les « Ce soir, essaie ceci ».
\item Lis \textbf{une fiche par soir}, pas plus. Chaque fiche tient sur une page, et son geste demande un quart d’heure environ.
\item Suis l’ordre la première fois, ou va directement à la fiche qui te parle. Si tu n’en lis qu’une, lis la fiche 2.
\item Ne cherche pas à tout appliquer. Un geste fait vraiment vaut mieux que quatorze gestes lus.
\end{itemize}

\textbf{Ce que ce livret n’est pas.} Ce n’est ni un cours de maths, ni une méthode miracle, ni un remplaçant de ton professeur. Il ne promet pas de bonnes notes. Il t’aide à lire ce que tes notes disent vraiment, et à choisir quoi faire ensuite.

\titre{Les quatorze fiches d’un coup d’œil}
\renewcommand{\arraystretch}{1.25}
\begin{tabularx}{\linewidth}{@{}r>{\bfseries}l X@{}}
\multicolumn{3}{@{}l}{\color{accent}\textbf{Partie I — Le verdict}}\\
1 & Le rang & Ce que ta note et ton rang disent, et ce qu’ils ne disent pas\\
2 & Les marches manquantes & Trouver la marche qui manque plutôt que de te dire « je suis nul »\\
3 & La main & Poser la question que tu gardes pour toi\\
4 & Le cahier du voisin & Te comparer sans te perdre\\
5 & La voix de dedans & Parler de ton échec avec précision\\[0.5em]
\multicolumn{3}{@{}l}{\color{accent}\textbf{Partie II — Le travail}}\\
6 & Redescendre & Réparer une base avec des exemples résolus\\
7 & Fermer le cahier & Se tester, revenir, étaler\\
8 & Dormir n’est pas tricher & Sommeil, ventre plein, téléphone loin\\
9 & L’encre rouge & Faire de tes erreurs un outil\\
10 & La copie double & Se préparer à l’examen, y compris à la peur\\[0.5em]
\multicolumn{3}{@{}l}{\color{accent}\textbf{Partie III — Le monde}}\\
11 & Ce qui ne dépend pas de toi & Faire la carte de ce que tu as\\
12 & Ce qu’on attend de toi & Parler d’orientation avec ta famille\\
13 & La lampe & Croiser deux sources pour une leçon difficile\\
14 & Le mot, encore & Déplacer le mot « nul » à sa juste place
\end{tabularx}

\partie{I}{Le verdict}

\fiche{1}{Le rang}{Chapitre 1 du livre}
\begin{aretenir}
Une note mesure \emph{une copie, un jour, sur une tâche}. Un rang dit \emph{où tu étais dans une file}, pas ce que tu vaux. Ni l’un ni l’autre ne dit « tu es comme ça ».
\end{aretenir}
\begin{itemize}
\item \textbf{Une note, c’est une mesure.} Comme la balance de la vendeuse qui pèse les tomates sans dire si elles sont bonnes : on s’en sert, en sachant ce qu’elle mesure. Dès les années 1930, Henri Piéron a montré que la même copie, confiée à plusieurs correcteurs sérieux, pouvait recevoir des notes très différentes. Ce que cela rend imprudent, ce n’est pas la note : c’est le \emph{verdict}.
\item \textbf{Le rang dépend de ta classe autant que de toi.} Les mêmes notes dans une autre classe donnent un autre rang. Il ressemble à une place dans une file d’attente : il dit qui est devant, qui est derrière, rien de la distance au guichet.
\item \textbf{La moyenne écrase.} Seize en maths et sept en français font onze et demi, et personne n’a jamais rendu une copie à onze et demi. Elle écrase aussi le temps : passer de quatre à huit, c’est doubler, mais la moyenne affiche six.
\item \textbf{Lis ton bulletin comme une carte, pas comme un jugement.} Un jugement se subit. Une carte sert à décider où aller la semaine prochaine.
\end{itemize}
\begin{cesoir}
\begin{enumerate}
\item Reprends ton dernier bulletin. À côté de chaque moyenne, écris \emph{combien de devoirs} elle résume, et de quel type ils étaient.
\item Choisis la matière où ta note est la plus basse. Cache-la, et écris la note que tu croyais avoir eue en sortant du devoir. Puis écris la vraie.
\item Regarde l’écart : il te dit si tu sais toi-même ce que tu sais. Dans trois jours, devine ta note avant de vérifier un exercice, et note les deux.
\end{enumerate}
\emph{Pour le cahier :} « Ma note de \_\_\_ dit que… » et « Elle ne dit pas que… ».
\end{cesoir}

\fiche{2}{Les marches manquantes}{Chapitre 2 du livre}
\begin{aretenir}
Le savoir scolaire est un escalier : chaque leçon s’appuie sur celles d’avant. Quand tu bloques, il manque souvent une marche plus bas. Remplace « je suis nul » par « il me manque une marche ».
\end{aretenir}
\begin{itemize}
\item \textbf{Une base absente se prend pour une incapacité.} Rodrigue, en troisième, bute sur $\frac{2}{3}+\frac{1}{4}$ : le calcul littéral du jour est compris, c’est l’addition de fractions qui manque. « Difficile » et « il me manque une marche » appellent des remèdes différents.
\item \textbf{Ta tête n’a que peu de place.} La mémoire de travail ne retient que quelques éléments à la fois. Ce qui est automatisé ne coûte presque rien ; une base non acquise occupe une main entière, et il n’en reste aucune pour la leçon du jour.
\item \textbf{Les marches manquantes se creusent.} Celui qui lit mal évite de lire, et l’écart grandit. C’est pourquoi mieux vaut réparer tôt.
\item \textbf{La honte garde les marches vides.} Réapprendre en troisième ce qu’on aurait dû savoir en CM2 est humiliant, alors on fait semblant. Personne n’est coupable : maladie, déménagement, professeur remplacé tard, leçon d’un jour d’absence.
\end{itemize}
\begin{cesoir}
\begin{enumerate}
\item Prends un exercice raté. Recopie la première étape où tu t’es arrêté.
\item Écris : « Pour faire ça, il fallait savoir… ». Puis : « Et pour savoir ça, il fallait savoir… ». Descends ligne par ligne, jusqu’à une chose dont tu es \emph{sûr}, au point de l’expliquer.
\item La marche juste au-dessus de cette chose sûre est la tienne. Entoure-la, note la date. \textbf{Ne la répare pas ce soir} : le diagnostic est déjà un travail.
\end{enumerate}
Si tu descends sans trouver de fond après cinq ou six lignes, garde la page : il te faudra quelqu’un pour descendre avec toi.
\end{cesoir}

\fiche{3}{La main}{Chapitre 3 du livre}
\begin{aretenir}
Ce qui te retient de poser une question, c’est le regard que tu imagines. Une question précise, posée par le bon canal, coûte peu et sert à beaucoup de monde.
\end{aretenir}
\begin{itemize}
\item \textbf{La timidité dépend de la situation.} Grâce chante alto devant une église pleine, mais n’ose pas lever la main : ce qu’elle craint, c’est d’être regardée au moment précis où elle ne sait pas.
\item \textbf{L’effet projecteur} (Gilovich, 2000) : on surestime ce que les autres remarquent de nous. Une question offre une pause, et on l’oublie dès la récréation.
\item \textbf{Tu n’es sans doute pas seul.} Le silence qui suit « C’est compris ? » est trompeur : chacun se croit le seul perdu. Poser la question rend service à la classe.
\item \textbf{La peur prend de la place.} L’inquiétude occupe une partie de l’espace où l’on raisonne. Une ligne non comprise ne reste pas au fond du cahier : elle revient, sous d’autres habits.
\item \textbf{Il y a d’autres chemins que la main levée} : la fin du cours, un papier laissé sur le bureau, le délégué, un camarade patient.
\end{itemize}
\begin{cesoir}
\begin{enumerate}
\item Prends un petit papier. Retrouve la dernière ligne de la leçon que tu comprends, puis la suivante, celle que tu as recopiée sans la suivre : ta question est entre les deux.
\item Écris-la précisément : « À la ligne 3, pourquoi le signe change ? » vaut mieux que « je n’ai pas compris ».
\item Donne-la au professeur à la fin du cours, ou laisse le papier sur son bureau. La semaine suivante, pose-en une à voix haute.
\end{enumerate}
\end{cesoir}
\begin{prudence}[Attention]
Moqueries quotidiennes, surnom qui humilie, messages qui te suivent chez toi : ce n’est plus de la timidité, c’est du \textbf{harcèlement}. Il ne se règle pas seul : dis-le à un adulte de confiance, et si le premier n’écoute pas, va voir le deuxième.
\end{prudence}

\fiche{4}{Le cahier du voisin}{Chapitre 4 du livre}
\begin{aretenir}
Te comparer est un réflexe normal. Le piège est de comparer ton brouillon à la copie au propre d’un autre. Pose la bonne question : pas « combien as-tu eu ? » mais \emph{« comment tu t’y prends, concrètement ? »}.
\end{aretenir}
\begin{itemize}
\item \textbf{On se compare à ce qui se voit} : la longueur d’une copie, une note dite dans la cour, une photo en blouse blanche. Ce qui se voit est vrai, mais ce n’est pas ce qui explique le reste. Ibrahim comparait une demi-page à six pages ; la juste comparaison aurait été une demi-page sans plan contre six pages précédées de quatre pages de brouillon.
\item \textbf{Une photo montre le résultat, jamais la fabrication.} Les écrans font cela à grande échelle : des vitrines, presque jamais d’arrière-boutiques.
\item \textbf{La mare compte.} Marsh et ses collègues ont observé qu’à niveau égal, on se juge moins bon dans une classe forte et meilleur dans une classe plus faible. Ce que tu crois valoir dépend en partie de la mare où l’on t’a mis.
\item \textbf{Une méthode se copie sans rien voler.} Ibrahim ne peut pas emprunter le père ni les trois années de dimanches, mais il peut emprunter le brouillon.
\item \textbf{La comparaison la plus honnête} est celle avec toi-même, quelques mois plus tôt. Encore faut-il avoir gardé une trace.
\end{itemize}
\begin{cesoir}
\begin{enumerate}
\item Choisis une personne accessible qui réussit dans la matière où tu peines, et écris son nom.
\item Demain, pose-lui la même question : « Comment tu t’y prends, concrètement, la veille d’un devoir ? » Si elle répond « je révise », demande ce qu’elle fait de ses mains.
\item Écoute sans comparer. Choisis \textbf{une seule habitude}, précise et gratuite, et tiens-la une semaine en notant les jours faits et oubliés.
\end{enumerate}
\end{cesoir}

\fiche{5}{La voix de dedans}{Chapitre 5 du livre}
\begin{aretenir}
L’histoire que tu te racontes après un échec pèse parfois autant que l’échec. « Je suis nul » juge sans avoir enquêté. Une phrase précise te donne une prise.
\end{aretenir}
\begin{itemize}
\item \textbf{Une explication peut te laisser une prise ou aucune.} Elle place la cause en toi ou hors de toi, passagère ou durable, sur un point précis ou sur tout. Selon l’histoire du soir, on ne fait pas la même chose le lendemain.
\item \textbf{À force, on cesse d’essayer.} Seligman et Maier ont décrit l’« impuissance apprise » ; la théorie a été révisée depuis (Seligman et Maier eux-mêmes, en 2016). Retiens l’idée prudente : répéter à quelqu’un que rien ne change rien le décourage.
\item \textbf{Parle-toi comme à un ami.} Des expériences (Kross et ses collègues, 2014) suggèrent que se parler à soi-même par son prénom ou par « tu » aide à prendre de la distance. L’effet est modeste et ne remplace rien, mais l’essai coûte une phrase.
\item \textbf{Le « pas encore ».} Pour ce qui compte à l’école (connaissances, méthodes), les capacités se travaillent. Mais la « mentalité de croissance » n’est pas une formule magique : les effets mesurés sont modestes et dépendent d’un entourage où l’effort est possible.
\end{itemize}
\begin{cesoir}
\begin{enumerate}
\item Écris, mot pour mot, la phrase que ta voix intérieure t’a dite aujourd’hui. Ne la commente pas.
\item Réécris-la précisément avec quatre questions : \emph{quoi, dans quelle matière, sur quel point, depuis quand} ? « Je suis nul en anglais » devient « Je confonds encore le prétérit et le \emph{present perfect}, depuis le début de l’année ».
\item Écris ce que tu dirais à ton meilleur ami avec la même copie, en lui disant « tu ». Relis les trois : la plus vraie est rarement la plus douce, et presque jamais la première.
\end{enumerate}
\end{cesoir}

\partie{II}{Le travail}

\fiche{6}{Redescendre}{Chapitre 6 du livre}
\begin{aretenir}
Pour réparer une base, ne cherche pas seul : étudie un exemple résolu, refais-le sans le regarder, et explique chaque ligne avec « parce que ». Puis remonte.
\end{aretenir}
\begin{itemize}
\item \textbf{Débutant, mieux vaut regarder faire.} Rodrigue a appris à ouvrir un téléphone en regardant tonton Fabrice, puis en faisant de plus en plus d’étapes. Sweller et la théorie de la charge cognitive décrivent la même chose : pour qui débute, étudier un exercice résolu est souvent plus efficace que chercher seul.
\item \textbf{Les « styles d’apprentissage » ne sont pas prouvés.} Une revue de Pashler et ses collègues (2008) n’a pas trouvé de preuve solide qu’adapter l’enseignement au style déclaré (visuel, auditif…) améliore l’apprentissage. Tu as des préférences, mais elles disent peu sur ta façon d’apprendre le mieux.
\item \textbf{Petit et régulier.} Peu d’exercices, chaque jour, jamais plus de vingt minutes : on garde l’envie de revenir.
\end{itemize}
\begin{cesoir}
\begin{enumerate}
\item Reprends la marche entourée (fiche 2). Trouve un livre d’une classe inférieure qui l’explique : celui d’un petit frère, d’un ancien, d’un voisin.
\item Lis un exemple résolu une fois, lentement. Cache-le avec une feuille et refais-le seul. Si tu bloques, découvre \emph{une seule} ligne, recache, continue.
\item Explique chaque étape à voix haute, en commençant par « parce que ». Si tu ne sais pas finir la phrase, tu viens de trouver où relire.
\item Pendant dix jours : trois exercices de cette marche par jour, vingt minutes au plus.
\item Quand tu réussis sans l’exemple, \textbf{remonte} : refais l’exercice de ta classe qui t’avait bloqué. On a réparé en bas ; on vérifie en haut.
\end{enumerate}
Dix jours réparent une marche. S’il en manque trois, compte un mois : c’est peu, à l’échelle d’une année.
\end{cesoir}

\fiche{7}{Fermer le cahier}{Chapitre 7 du livre}
\begin{aretenir}
Relire donne le sentiment de savoir. \textbf{Se tester}, puis \textbf{revenir} à intervalles de plus en plus longs, fait réellement retenir.
\end{aretenir}
\begin{itemize}
\item \textbf{L’illusion de maîtrise.} Une page relue quatre fois paraît familière ; une solution lue avant d’avoir cherché paraît facile. Roediger et Karpicke (2006) : cinq minutes après, ceux qui avaient relu s’en sortaient mieux ; une semaine plus tard, ceux qui s’étaient testés retenaient nettement plus, et les relecteurs se croyaient mieux préparés.
\item \textbf{Ce qui marche le mieux.} La revue de Dunlosky et ses collègues (2013) place en tête se tester et espacer ses révisions. Relire et surligner sont parmi les moins utiles. Garde tes surligneurs si tu les aimes : sers-t’en pour marquer ce sur quoi tu vas t’interroger.
\item \textbf{Oublier est normal.} Depuis Ebbinghaus, on sait que l’oubli est le fonctionnement ordinaire de toute mémoire. On ne l’arrête pas ; on le ralentit en \emph{revenant} : le lendemain, quelques jours plus tard, la semaine suivante.
\item \textbf{Les « difficultés désirables »} (Bjork) : l’effort de faire revenir ce qu’on a appris aide ; ce qui coule tout seul s’efface vite. Mélanger les types d’exercices aide aussi à reconnaître de quel problème il s’agit.
\end{itemize}
\begin{cesoir}
\begin{enumerate}
\item Après une leçon, ferme le cahier, pose-le loin. Écris pendant dix minutes tout ce dont tu te souviens, sans regarder. Marque d’un « ? » ce dont tu doutes.
\item Rouvre et corrige avec un autre stylo : ces corrections dessinent ce qu’il faut retravailler.
\item Écris trois questions qui t’obligent à produire (« Pourquoi la deuxième étape vient-elle avant la troisième ? »), avec des dates : demain, dans trois ou quatre jours, la semaine suivante. Réponds sans notes, vérifie ensuite.
\end{enumerate}
\textbf{La boîte de Leitner} (sans courant ni connexion). Un carton à trois compartiments, des cartes question/réponse. Compartiment 1 chaque jour, compartiment 2 tous les deux ou trois jours, compartiment 3 chaque semaine. Bonne réponse : la carte avance. Erreur : retour au compartiment 1.
\end{cesoir}

\fiche{8}{Dormir n’est pas tricher}{Chapitre 8 du livre}
\begin{aretenir}
Dormir fait partie de la révision. Un ventre vide, un téléphone sur la table et des nuits courtes coûtent de l’attention, sans que tu t’en rendes compte.
\end{aretenir}
\begin{itemize}
\item \textbf{La nuit travaille.} Les travaux de Jan Born et d’autres équipes suggèrent qu’une partie de ce qu’on a appris dans la journée est reprise et stabilisée pendant le sommeil. On ne sait pas encore bien comment : méfie-toi des affirmations spectaculaires.
\item \textbf{À l’adolescence, l’horloge se décale} (Carskadon) : l’envie de dormir arrive plus tard. Ce n’est pas de la paresse.
\item \textbf{Le corps entre en classe avec la tête.} Un ventre vide à 10~h pense mal ; manger quelque chose, même peu, est un geste de travail. L’activité physique a des effets favorables mais \emph{modestes} sur l’humeur et l’attention : méfie-toi de ce qui se vend en flacon.
\item \textbf{Changer de tâche a un coût.} « Faire deux choses à la fois » est une alternance rapide, et chaque aller-retour se paie. Une notification t’éloigne du début de ta phrase. Une étude (Ward, 2017) trouvait même un effet du téléphone posé sur la table ; il n’a pas toujours été retrouvé depuis. Le pari reste sans risque : éloigner le téléphone ne coûte rien.
\end{itemize}
\begin{cesoir}
\begin{enumerate}
\item Repère ton prochain devoir important. Compte sept jours en arrière et choisis \emph{une} heure de coucher, réaliste pour ta maison. Tiens-la surtout les soirs où tu n’as pas fini ; note le reste en tête du lendemain.
\item La veille : ne commence rien de nouveau après le dîner. Reprends les questions que tu t’es écrites, réponds sans regarder, vérifie, éteins.
\item Pour tes séances de travail : confie ton téléphone à quelqu’un pour 45 minutes. Ensuite, écris en une ligne ce que tu as \emph{réellement} fait.
\end{enumerate}
\end{cesoir}

\fiche{9}{L’encre rouge}{Chapitre 9 du livre}
\begin{aretenir}
Une erreur est une information. Pour qu’elle serve, il faut la nommer, la comprendre, puis la refaire plus tard sans regarder.
\end{aretenir}
\begin{itemize}
\item \textbf{Les erreurs faites avec assurance se retiennent bien une fois corrigées} (Metcalfe, effet d’hypercorrection) : la surprise réveille l’attention.
\item \textbf{Un bon retour d’information} dit où tu en es, où tu devais arriver, et ce qu’il faut faire ensuite (Hattie et Timperley). Une note seule ne dit que le premier point : c’est à toi d’écrire les deux autres.
\item \textbf{Chercher avant de voir la solution} peut aider, dans certaines conditions ; dans d’autres, cela décourage seulement. Cinq minutes suffisent pour que la question devienne la tienne.
\item \textbf{Expliquer à d’autres} oblige à mettre de l’ordre, et fait découvrir les trous. Le tutorat entre élèves profite souvent aux deux.
\end{itemize}
\begin{cesoir}
\begin{enumerate}
\item Prends ta dernière copie corrigée. Dans un carnet d’erreurs (ou à l’envers à la fin du brouillon), pour \emph{chaque} erreur écris quatre phrases : \emph{ce que j’ai fait} ; \emph{pourquoi} ; \emph{ce qu’il fallait faire} (règle et raison) ; \emph{un exercice semblable} avec une date, dans trois jours.
\item Ce jour-là, refais l’exercice sans regarder. Réussi : reviens une semaine plus tard. Raté : retour à trois jours.
\item Note la \textbf{famille} de chaque erreur : lecture, attention, méthode ou notion. Si la même famille revient trois fois, tu as trouvé ton point faible du moment.
\item Cette semaine, propose à deux camarades une heure de « tour de tableau » : chacun explique un exercice sans notes, et les autres demandent « pourquoi ? » à chaque étape.
\end{enumerate}
\end{cesoir}

\fiche{10}{La copie double}{Chapitre 10 du livre}
\begin{aretenir}
La meilleure préparation à un examen, c’est l’examen : une épreuve entière, dans les vraies conditions. Et trois gestes écrits à l’avance pour les dix premières minutes.
\end{aretenir}
\begin{itemize}
\item \textbf{Fais une annale en vrai} : table dégagée, pas de téléphone, chronomètre, copie double, sans ouvrir le cours. Tu apprends des choses qu’aucune leçon n’enseigne : le temps passe plus vite quand il est compté, la main fatigue, on bâcle ce qu’on craint.
\item \textbf{Le jour J, garde un ordre simple.} Respirer lentement (l’expiration plus longue que l’inspiration). Lire \emph{tout} le sujet, marquer au crayon ce que tu sais et ce qui te fait peur. Commencer par ce que tu sais, pour faire baisser le bruit dans ta tête. Répartir le temps selon les points, et quitter chaque exercice à l’heure prévue.
\item \textbf{Lire le stress autrement.} Jamieson et ses collègues (2010) ont observé que lire que le cœur rapide peut être la préparation du corps à l’effort aidait des étudiants en maths. Une seule étude ne fait pas une loi : c’est une façon de lire les signes, pas une recette.
\item \textbf{À ne pas retenir comme une preuve.} Écrire ses inquiétudes avant l’examen (Ramirez et Beilock, 2011) n’a pas été retrouvé lors d’une grande reprise en 2018.
\end{itemize}
\begin{cesoir}
\begin{enumerate}
\item Écris dès ce soir \textbf{tes trois gestes} des dix premières minutes, chacun avec un verbe au début (ex. : respirer, lire tout le sujet, choisir l’exercice par lequel commencer). Recopie-les sur un ticket glissé dans ta trousse.
\item Réserve, dans les deux ou trois semaines à venir, un créneau entier, à l’heure exacte de l’épreuve que tu redoutes, pour une annale avec corrigé.
\item Le lendemain, corrige-toi comme un correcteur sévère. Pour chaque point perdu : notion, consigne mal lue, calcul ou temps ? Ce relevé en dit plus que la note.
\end{enumerate}
\end{cesoir}

\partie{III}{Le monde}

\fiche{11}{Ce qui ne dépend pas de toi}{Chapitre 11 du livre}
\begin{aretenir}
Tout n’est pas question de volonté : le courant, les livres, le calme, le coût des annales ne dépendent pas de toi. Mais ce qui dépend un peu de toi est souvent plus grand que tu ne le crois. Commence par en faire la carte.
\end{aretenir}
\begin{itemize}
\item \textbf{Ce n’est pas ta faute.} Bourdieu et Passeron (\emph{Les Héritiers}, 1964) ont montré que l’école prend souvent pour des qualités personnelles des avantages reçus en naissant, et les récompense comme du mérite.
\item \textbf{Le manque pèse sur l’attention.} Mullainathan et Shafir suggèrent que les soucis d’argent occupent l’esprit, même quand on pense à autre chose (certains résultats sont discutés). Une inquiétude écrite avec une date a moins besoin d’être répétée.
\item \textbf{La distinction d’Épictète}, il y a presque deux mille ans : certaines choses dépendent de nous (nos jugements, nos élans), d’autres non (le corps, les biens, la réputation).
\item \textbf{La frontière se déplace quand on s’associe} : les annales d’une amie, une question écrite à trois, une bibliothèque à l’heure calme.
\end{itemize}
\begin{cesoir}
\begin{enumerate}
\item Sur une page, dresse la carte de ce que tu \emph{as} : les \textbf{heures} (même vingt minutes, le trajet en car), les \textbf{endroits} (bibliothèque, salle vide, cour d’un voisin) avec l’heure, les \textbf{personnes} qui savent quelque chose (avec ce qu’elles savent), les \textbf{objets} (un manuel même ancien, des annales, une lampe).
\item Choisis \textbf{un seul levier} et utilise-le cette semaine : demander le cahier de l’ancien, aller à la bibliothèque à l’heure calme, poser ta question au camarade.
\item Dans sept jours, note ce qu’il t’a donné, même si c’est peu.
\end{enumerate}
\end{cesoir}

\fiche{12}{Ce qu’on attend de toi}{Chapitre 12 du livre}
\begin{aretenir}
On choisit souvent un métier sur une image : une blouse, un bureau, une boutique. Une image ne dit rien d’une journée. Pour choisir, ou discuter un choix, il faut parler à ceux qui exercent vraiment le métier. Et commencer, avec ta famille, par écouter leur peur.
\end{aretenir}
\begin{itemize}
\item \textbf{S’engager durablement} suppose de se sentir compétent, d’avoir une part de choix et de se sentir relié aux autres (Deci et Ryan, théorie de l’autodétermination ; les détails sont débattus, l’idée est solide).
\item \textbf{La conversation avance mieux} avec un projet concret (quel établissement, quelle série, ce qu’on y apprend, ce que font ceux qui en sortent, le plan B) qu’avec un désir. Il faut accepter de ne pas tout obtenir le premier soir ; proposer d’aller voir ensemble aide.
\end{itemize}
\begin{cesoir}
\begin{enumerate}
\item Choisis un adulte qui fait le métier qui t’intéresse (ou celui que tes parents imaginent pour toi). Écris, en une phrase, le métier \emph{tel que tu l’imagines}.
\item Demande-lui dix minutes et pose quatre questions : Qu’est-ce que vous faites vraiment dans une journée ordinaire ? Qu’avez-vous dû apprendre, à l’école et après ? Qu’est-ce qui est dur, que les gens ne voient pas ? Qu’est-ce que vous n’aviez pas imaginé avant de commencer ?
\item Le soir même, note ses réponses avec ses mots. Souligne ce qui était juste dans ta phrase, barre le reste.
\item Prépare \textbf{une phrase d’ouverture} pour tes parents, qui commence par \emph{leur} inquiétude : « Je sais que tu as peur que… J’ai demandé à quelqu’un qui fait ce métier comment ça se passe vraiment, et je voudrais te dire ce qu’il m’a répondu. »
\end{enumerate}
\end{cesoir}

\fiche{13}{La lampe}{Chapitre 13 du livre}
\begin{aretenir}
Une lampe éclaire la marche suivante ; elle ne monte pas l’escalier à ta place. Pour une leçon difficile, croise \textbf{deux} sources de nature différente, et compare-les.
\end{aretenir}
\begin{itemize}
\item \textbf{Une source seule peut te tromper}, un manuel comme une application. Deux sources qui disent la même chose te donnent une base plus solide ; qui se contredisent, te donnent une \emph{question pour ton professeur}.
\item \textbf{Aucune lampe ne remplace la classe}, le professeur qui te voit, ni un camarade patient. Les outils sur téléphone supposent de la batterie et un forfait, et ceux qui manquent le plus de cours peuvent manquer aussi de cela. Une lampe se partage : un téléphone peut éclairer trois cahiers, et une méthode se recopie pour le jour où la batterie lâche.
\item \textbf{OnBuch en fait partie, parmi d’autres} : des cours structurés leçon par leçon, des fiches de travaux dirigés à trois niveaux, des corrigés, des fiches méthode, construits sur les programmes du MINESEC. Il peut y rester des erreurs : c’est une raison de plus de croiser les sources.
\end{itemize}
\begin{cesoir}
\begin{enumerate}
\item Prends la leçon qui te résiste le plus et choisis \textbf{deux lampes} nommées précisément (quel manuel, quelle personne), de nature différente. Écris leurs noms et la date.
\item Demain, lis ou écoute la première et résume en quelques lignes avec tes mots. Après-demain, la seconde.
\item Compare les deux résumés ligne par ligne. Elles se contredisent ? Écris ta question pour le professeur en nommant l’endroit exact du désaccord.
\item Dans trois jours, sans rien relire, refais un exercice de cette leçon : tu sauras si les lampes ont éclairé quelque chose en toi, ou seulement ta page.
\end{enumerate}
\end{cesoir}

\fiche{14}{Le mot, encore}{Épilogue du livre}
\begin{aretenir}
On ne guérit pas d’un mot, mais on peut le \emph{déplacer} : le décoller du nom où quelqu’un l’a posé et le remettre là où il a un sens, sur un résultat, sur une copie, sur une marche.
\end{aretenir}
\begin{itemize}
\item \textbf{« Nul », en maths, ça veut dire zéro.} Dans $7\times5\times0$, il suffit d’un seul zéro pour que le résultat ait l’air nul, mais le sept est toujours là, et le cinq aussi. Un zéro dans un produit ne rend pas les autres facteurs nuls.
\item \textbf{Les zéros se cherchent.} Un quatre en cinquième, un « je suis nulle » à dix ans devant un cahier de vacances ne prouvent pas ce que tu crois. Ils disent qu’un facteur manque : une base, une méthode, du temps, de la lumière. Ces facteurs-là se trouvent, et se corrigent.
\item \textbf{Ce que tu apprends, tu peux le transmettre.} Rodrigue explique les fractions à sa cousine avec des arachides ; c’est en expliquant qu’on s’aperçoit de ce qu’on sait.
\end{itemize}
\begin{cesoir}
\begin{enumerate}
\item Retourne à la première page de ton cahier de brouillon. Relis la phrase que quelqu’un a dite sur toi. \textbf{Ne la rature pas.}
\item En dessous, réécris-la à sa juste place : ce qu’elle disait d’un moment, d’une copie, d’une marche. Pas de toi.
\end{enumerate}
\end{cesoir}

\titre{Ta première semaine}
Si tu n’as qu’une semaine, voici un parcours léger. Un quart d’heure par soir suffit.
\renewcommand{\arraystretch}{1.35}
\begin{tabularx}{\linewidth}{@{}l X@{}}
\toprule
\textbf{Jour} & \textbf{Geste}\\\midrule
Lundi & \textbf{Fiche 1.} Note devinée contre note réelle, dans la matière la plus basse.\\
Mardi & \textbf{Fiche 2.} Descends jusqu’à la marche qui manque. Entoure-la, ne la répare pas encore.\\
Mercredi & \textbf{Fiche 3.} Écris ta question précise, donne-la au professeur.\\
Jeudi & \textbf{Fiche 5.} Ta phrase intérieure, réécrite en trois lignes.\\
Vendredi & \textbf{Fiche 6.} Premier exemple résolu, refait sans regarder, avec des « parce que ».\\
Samedi & \textbf{Fiche 7.} Dix minutes de cahier fermé, trois questions datées.\\
Dimanche & \textbf{Fiche 8.} Choisis ton heure de coucher pour la semaine qui vient.\\\bottomrule
\end{tabularx}

\vspace{0.8em}
Ensuite, avance à ton rythme : la fiche 9 dès que tu reçois une copie corrigée, la 10 trois semaines avant un examen, les fiches 4, 11, 12 et 13 quand le sujet se présente, la 14 quand tu veux refermer la boucle.

\titre{Ce qui est solide, ce qui l’est moins}
Ce livret, comme le livre, essaie de dire à chaque fois si un résultat est solide, discuté ou fragile. La recherche avance et se corrige.
\renewcommand{\arraystretch}{1.3}
\begin{tabularx}{\linewidth}{@{}>{\bfseries}p{2.6cm} X@{}}
\toprule
Solide & Se tester plutôt que relire ; espacer les révisions ; étudier des exemples résolus quand on débute ; recevoir un retour d’information utile.\\\midrule
Plausible, effets modestes & Dormir assez ; manger ; bouger ; éloigner le téléphone ; se parler à la deuxième personne ; lire son stress comme une préparation du corps.\\\midrule
Discuté & La mentalité de croissance (effets modestes en moyenne) ; l’impuissance apprise (théorie révisée) ; l’effet Pygmalion (réel mais limité).\\\midrule
Non soutenu ou non retrouvé & Les « styles d’apprentissage » ; l’effet de l’écriture des inquiétudes avant un examen (non retrouvé en 2018).\\\bottomrule
\end{tabularx}

\vspace{0.8em}
Les références complètes se trouvent dans la note de l’auteur du livre. Aucune n’est nécessaire pour faire les gestes de ce livret : un cahier de brouillon suffit.

\titre{Si ça ne va pas}
Cette page n’a rien à t’apprendre sur les notes.

Il arrive que la fatigue ou la tristesse dure plus longtemps qu’une mauvaise semaine : le sommeil ne vient plus, tu ne manges plus vraiment, ce que tu aimais ne te fait plus rien, ta voix de dedans ne s’arrête plus. Il arrive aussi qu’une idée vienne, celle de disparaître ou de te faire du mal. Ou que quelqu’un te frappe, te menace, t’humilie.

Si c’est ton cas, la méthode ne suffit plus, et tu n’as pas à t’en sortir sans aide. Parles-en à quelqu’un, aujourd’hui si tu peux : un parent, un frère ou une sœur aînée, un professeur, le surveillant général, l’infirmière ou le conseiller d’orientation de ton établissement, un médecin ou un centre de santé. Tu n’as pas besoin de trouver les bons mots : « Ça ne va pas, et je ne sais pas quoi faire » suffit. Si la première personne ne comprend pas, va voir la suivante.

Si le danger est immédiat, fais-toi accompagner tout de suite : réveille quelqu’un dans la maison, ou va vers le centre de santé le plus proche.

Tu peux aussi joindre l’auteur : +237 678 43 85 57, 689 72 19 23 ou ludovic@onbuch.site. Ce n’est ni un service d’urgence ni un service de soin, mais si tu ne sais pas vers qui aller, écris ou appelle : il te mettra en contact avec quelqu’un de qualifié.

\titre{Pour les parents et les enseignants}
Vous n’avez pas besoin de connaître la leçon pour aider.
\begin{aretenir}[Une question à poser]
Plutôt que « Tu as appris ta leçon ? », qui appelle un oui sans valeur (l’enfant a souvent relu et croit savoir), essayez : \textbf{« Explique-moi ce que tu as compris. »} Vous entendrez très bien si l’explication tient debout.
\end{aretenir}
Quelques gestes du livret se font mieux à deux : garder le téléphone pendant quarante-cinq minutes (fiche 8), poser les questions de l’annale en conditions réelles (fiche 10), écouter la peur avant de discuter l’orientation (fiche 12). Si un enfant change sans bruit, s’isole, ne dort plus ou parle de lui avec une dureté nouvelle, ne le prenez pas pour de la paresse : arrêtez-vous, posez la question franchement et écoutez sans juger (voir « Si ça ne va pas »).

\titre{Et le livre ?}
Ce livret est un avant-goût. Il garde les gestes et les raisons. Il laisse de côté ce qui fait le livre : Rodrigue, Grâce, Divine, Ibrahim, Nadège et Aïcha ; monsieur Kamga et sa lampe qui s’éteint ; la tontine de Mama Odile ; une lettre à Aïcha ; les nuances et les sources. Ces personnes sont composées à partir de situations réelles et typiques ; aucune n’est une personne précise.

Si une fiche t’a parlé, le chapitre correspondant te dira pourquoi elle marche, et ce qu’elle ne peut pas faire.

\vspace{1em}
Pour écrire à l’auteur à propos du livre ou d’OnBuch : +237 678 43 85 57, 689 72 19 23 ou ludovic@onbuch.site.
\end{document}
